% Four optional replacement/addition paragraphs. Choose by argumentative value,
% not to meet a citation quota. Metadata/access depth: source-register.json.

% Direct domain comparator; place after the prospective 2026 paragraph.
A separate archival study compares channeling-related and ufological material through an additional computational layer: both collections are queried using the same UFODex system, and embeddings compare its generated answers to ten matched questions \citep{wahbeh2026records}. Its comparison units are therefore AI-mediated responses, whereas ours are representations of published titles and available captions. The resulting cosine values cannot be pooled, and resemblance in either design does not establish an independent source of information.

% Anthropology context; these two sources were available at abstract level only.
Anthropological accounts also locate channeling and mediumship in learned practices and interactions. \citet{klinoron2014channel} describes participation in an Israeli channeling course, while \citet{tomlinson2019medium} examines voice, agency, and audience response in Australian Spiritualism. These different settings motivate attention to performance and context rather than isolated utterances. We use their publisher abstracts for this limited contextual framing; they do not support classifying our indexed practitioners or inferring the origin of their communications.

% Highest-priority platform-method addition; Tube2Vec is an ICWSM Poster Paper.
Tube2Vec is a close computational precedent: its content representation encodes titles and descriptions with all-MiniLM-L6-v2, sums their vectors within each video, and averages the 50 most recent videos per account \citep{boesinger2024tube2vec}. It compares these vectors with social-sharing and recommendation representations and uses human judgments for task-specific evaluation. Its filtered sample emphasizes large, mostly English accounts. Our contribution is not the averaging operation, but the comparison of matched text layers and the consequences of provenance and sampling decisions in a small, selected corpus; its human validation does not transfer to ours.

% Replace/extend the current religion/small-community bridge; this is not ET evidence.
The religious-video analysis of \citet{chen2026digital} combines narrative, visual and comment coding with human validation, while \citet{claisse2023faith} studies twelve Buddhist community members through interviews, reflexive thematic analysis and first-author autoethnography. Together these studies show two different ways of linking mediated communication to context. Our public-record analysis has neither their participant accounts nor their human-validated labels, and therefore makes narrower claims about observable metadata and text representations.
